% MAIN SNAPSHOT -- LSPV/RMV::COCYCLEPOTENTIAL1::2026-10-15T13:59:00+09:00
% LSPV-5 REPORT::COCYCLEPOTENTIAL1::2026-10-15T13:59:00+09:00
\documentclass{article}
\usepackage{amsmath,amssymb,amsthm,hyperref}
\newtheorem{theorem}{Theorem}
\title{COCYCLEPOTENTIAL1: Source-Checked Nonzero Memory Parity Changes}
\author{LSPV-5}
\date{15 October 2026}
\begin{document}\maketitle
\noindent\textbf{Status:} Restricted theorem PROVED/EXACT, review pending; P versus NP OPEN. This is a condensed publication. A separate full ten-page TeX/PDF and six reproducibility checkpoints preserve the expanded proofs and tests.

\section{IDEAS PROPOSED BEFORE LITERATURE}
The pre-search proposals were separately frozen: \textbf{CANDIDATE A}, compensate fixed protected-write parity by an exact finite-control potential (falsifier: inconsistent source paths); \textbf{CANDIDATE B}, use multiple bank parities while allowing data-dependent addresses but not runtime-selected bank labels (falsifier: variable bank); \textbf{CANDIDATE C}, derive a canonical two-output set from one source circuit (falsifier: forged third word); \textbf{CANDIDATE D}, audit same-oracle and completion-time WRAPPER transfers (falsifier: oracle-private code preserving the invariant). Each proposal included its missing instruction-typing hypothesis; later revisions are stated separately below.

\section{LITERATURE CHECK AND REVISIONS}
Finite potentials are established signed-graph balance mathematics (Zaslavsky, 1981). The checked source grammar is consistent with the proof-carrying code distinction between typed execution and untrusted annotation (Necula, 1997). The original Hirahara theorem 3.9 gives P-immunity of superpolynomial time-bounded $K^t$ randomness, not Levin $Kt$. Baker--Gill--Solovay (1975) gives an oracle with $P^O=NP^O$; only its original journal abstract and bibliography, not the full paper, were directly available. Aaronson--Wigderson (2008) requires a separate algebrization test. Candidates A--C survive with enforced typing; D fails as a new nonrelativization proof. Cryptographic hardness is not assumed.

\section{CONSTRUCTION OR PROOF ATTEMPT}
A source of length $N$ explicitly defines a finite control multigraph $(Q,E)$, $k\le N$ protected arrays of length $2^d$, $d\le N$, and on every edge an explicit list of XOR toggles, each containing a fixed bank number but an arbitrary dynamically computed address. Private code can read protected memory and perform arbitrary calculations, branches and delays. The actual typed interpreter forbids raw SET, runtime-selected bank identity and untyped output. Initial arrays have an explicit compact all-zero grammar; the verifier does not enumerate them.

Let $\sigma(X)\in\mathbb F_2^k$ be the vector of parity of all cells in each bank, and let $\gamma_e\in\mathbb F_2^k$ be the XOR of the unit bank vectors for all toggles on edge $e$. Labels are recomputed from source code.

\begin{theorem}[Exact source potential]
There exists $\phi:Q\to\mathbb F_2^k$ such that
\[
\phi(u)\oplus\phi(v)=\gamma_e\qquad(e:u\to v)
\]
if and only if \emph{every undirected cycle} of the labelled multigraph, including parallel edges and self-loops, has XOR-label sum zero. A canonical solution is polynomially computable.
\end{theorem}
\begin{proof}
Necessity follows by telescoping endpoint differences on undirected closed walks. For sufficiency, fix zero potential at each spanning-forest root and propagate it through tree edges. Every remaining edge is consistent exactly when its fundamental undirected cycle has zero label. This also checks loops and parallel edges.
\end{proof}

\begin{theorem}[Long computation and global binding]
At each legal edge-boundary state $(q,X)$, the quantity
\[
Z=\sigma(X)\oplus\phi(q)
\]
is constant. Inside an edge from $u$, after $r$ toggles, the stronger invariant is
\[
Z=\sigma(X)\oplus\phi(u)\oplus\Delta_e(r),
\quad \Delta_e(r)=\bigoplus_{i\le r}\mathbf e_{j_i}.
\]
A verifier which accepts only the two values $G(Z,0),G(Z,1)$ for one source-owned output circuit $G$ globally satisfies
\[
|\{y:\exists\pi\ V(S,y,\pi)=1\}|\le2.
\]
The one verifier has conservative bounds $T=O(N^4)$, $S=O(N^2)$, $|\pi|=O(N^2)$. These bounds are for certificate checking, not for underlying execution or SAT decision.
\end{theorem}
\begin{proof}
An XOR toggle changes the parity of precisely its statically designated bank, regardless of address collisions. The corresponding prefix vector changes by the same amount, preserving the expression even between ordinary Turing machine steps. At edge completion $\Delta_e=\gamma_e=\phi(u)\oplus\phi(v)$ gives the boundary invariant. Source-checking the graph canonically fixes $Z$ from initial memory. The verifier calculates the only permitted output pair directly from $G$, independently of the proof string; hence the two-output bound holds \emph{over all accepted proofs}. Parsing, graph checking and circuit evaluation have the stated conservative polynomial resources.
\end{proof}

\paragraph{Smallest falsification.}
Two distinct edges $0\to1$ labelled $0$ and $1$ form a directed acyclic multigraph with inconsistent potential. A simple directed acyclic triangle $(0\to1,0),(1\to2,0),(0\to2,1)$ is also inconsistent. Therefore tests of only \emph{directed} cycles are FALSIFIED/RETIRED. Untyped SET and runtime-selected bank operations also falsify generalizing the invariant to arbitrary programs. Source enforcement, not annotations, is essential.

\paragraph{Large configuration control.}
A short typed source can toggle each of $2^d$ physically stored bits, leaving exponentially wide intermediate memory while maintaining the invariant. The verifier never scans the memory. The actual execution is exponential and is not accelerated by the certificate.

\paragraph{Matched-oracle and WRAPPER.}
For the same oracle $O$ supplied to all machines, polynomial suffix enumeration proves $R_{K_U^{t,O}}$ is $P^O$-immune for superpolynomial time bound $t$. If $P^O=NP^O$, a universal polynomially verifiable two-output certificate covering all time-bounded outputs of programs $|p|\le n-2$ would define a $P^O$ set $E_n$ of at most $2^n-2$ strings, whose infinite complement would be a $P^O$ subset of the immune language. Contradiction. This retains the earlier PVG-4-0001 diagnostic; it does not construct universal certificates. The positive source-phase theorem itself relativizes. A bitstream prefix revealed at completion can be an offline input after release, regardless of online unpredictability. Nonalgebrization OPEN.

\section{TESTS AND OUTCOME}
Six independently checkpointed runs, seed 1510152026, checked 65,349 finite cases, 429,657 assertions, zero detected discrepancies. Exhaustive graph-potential cases, randomized microsteps, conflicting paths, untyped hostile controls, memory sweeps and forged outputs were included. Physical sweeps reached $2^{13}$ cells; $2^{40}$, $2^{64}$ and $2^{128}$ arrays were symbolic only. Evidence is COMPUTATIONAL, not asymptotic verification.

\section{FIVE-TRACK STATUS AND NEXT TARGET}
LSPV-1: restricted hereditary contraction, arbitrary all-input contractor OPEN. LSPV-2: compact global certificates despite exponential LP trees, uniform arbitrary-decider proof transport OPEN and RS5 (H) CONDITIONAL. LSPV-3: small quotient-vertex-cover ternary CSP algorithm, satisfiable UWI common-input obstruction OPEN. LSPV-4: fixed-address equality-only extension evaluation, unrestricted same-machine nonrelativizing simulation OPEN. LSPV-5: new typed nonzero-bank-parity certificate, universal machine transport OPEN. Next target: a source-checkable nontrivial class of data-dependent protected signature changes without runtime guards.

\begin{thebibliography}{5}
\bibitem{z} Zaslavsky, Characterizations of signed graphs, J. Graph Theory, 1981.
\bibitem{n} Necula, Proof-Carrying Code, POPL, 1997.
\bibitem{h} Hirahara, Unexpected Hardness Results for Kolmogorov Complexity Under Uniform Reductions, STOC 2020, \url{https://eccc.weizmann.ac.il/report/2020/050/download}.
\bibitem{b} Baker, Gill, Solovay, Relativizations of the P = ? NP Question, SIAM J. Comput., 1975.
\bibitem{a} Aaronson, Wigderson, Algebrization: A New Barrier in Complexity Theory, STOC 2008.
\end{thebibliography}
\end{document}
% FINAL_VERIFICATION_SENTINEL::LSPV5::COCYCLEPOTENTIAL1::2026-10-15_1359::P_VS_NP_OPEN
% END_OF_VERIFIED_MAIN::COCYCLEPOTENTIAL1